% Plastic displacement of a faceted grain boundary:
% facet construction (GbFacet) and evaluation of the displacement field
% on and away from the boundary (GbFacet, GbContinuum).
%
% Self-contained: compiles with pdflatex, needs only amsmath/amssymb.
% To lift into a larger document, take everything between \begin{document}
% and \end{document} and drop the \section* levels one step.

\documentclass[11pt]{article}
\usepackage{amsmath}
\usepackage{amssymb}
\usepackage[margin=1in]{geometry}

\newcommand{\Lat}[1]{\mathcal{#1}}
\newcommand{\A}{\Lat{A}}
\newcommand{\B}{\Lat{B}}
\newcommand{\nhat}{\hat{\mathbf{n}}}
\newcommand{\ubar}{\bar{\mathbf{u}}}
\newcommand{\Sinf}{S_{\infty}}

\title{Plastic displacement of a faceted grain boundary:\\
       facet construction and evaluation}
\date{}

\begin{document}
\maketitle

\section*{Setting}

A mesostate is specified by a set of engaged translation--shift pairs
$(\mathbf{t}_i, \mathbf{s}_i)$. Each pair contributes one node to each grain,
\begin{equation}
  \mathbf{x}^{\A}_i = \mathbf{s}_i - \tfrac{1}{2}\mathbf{t}_i \in \A,
  \qquad
  \mathbf{x}^{\B}_i = \mathbf{s}_i + \tfrac{1}{2}\mathbf{t}_i \in \B,
  \label{eq:nodes}
\end{equation}
carrying the nodal displacements
\begin{equation}
  \mathbf{u}^{\A}_i = +\tfrac{1}{2}\mathbf{t}_i,
  \qquad
  \mathbf{u}^{\B}_i = -\tfrac{1}{2}\mathbf{t}_i ,
  \label{eq:nodal-displacements}
\end{equation}
so that both nodes deform onto the same point,
\begin{equation}
  \mathbf{x}^{\A}_i + \mathbf{u}^{\A}_i
  = \mathbf{x}^{\B}_i + \mathbf{u}^{\B}_i
  = \mathbf{s}_i .
  \label{eq:coincidence}
\end{equation}
The boundary is doubly periodic, with in-plane period vectors
$\mathbf{p}_1, \mathbf{p}_2$ spanning the plane whose unit normal is $\nhat$.
Write $\mathbf{T}_{mn} = m\mathbf{p}_1 + n\mathbf{p}_2$ for a period
translation.

Each grain is bounded by a triangulated surface, its \emph{facet}, carrying the
nodal displacement field of~\eqref{eq:nodal-displacements}. The plastic
displacement is the field generated by that surface treated as a
displacement-jump sheet.

\section*{Facet construction}

\subsection*{Periodic triangulation on the torus}

Let $\mathbf{e}_1, \mathbf{e}_2$ complete an orthonormal in-plane frame with
$\mathbf{e}_1 \times \mathbf{e}_2 = \nhat$, and let
\begin{equation}
  \mathsf{M} =
  \begin{bmatrix}
    \mathbf{p}_1\cdot\mathbf{e}_1 & \mathbf{p}_2\cdot\mathbf{e}_1 \\
    \mathbf{p}_1\cdot\mathbf{e}_2 & \mathbf{p}_2\cdot\mathbf{e}_2
  \end{bmatrix}
  \label{eq:torus-map}
\end{equation}
map the period parallelogram onto the unit-square torus $[0,1)^2$. Each node is
projected and reduced to its canonical torus coordinate,
\begin{equation}
  \mathbf{c}_i = \mathsf{M}^{-1}
  \begin{bmatrix} \mathbf{x}_i\cdot\mathbf{e}_1 \\
                  \mathbf{x}_i\cdot\mathbf{e}_2 \end{bmatrix},
  \qquad
  \boldsymbol{\xi}_i = \mathbf{c}_i - \lfloor \mathbf{c}_i \rfloor
  \; \in [0,1)^2 ,
  \label{eq:canonical-uv}
\end{equation}
the discarded integer part $-\lfloor \mathbf{c}_i \rfloor$ being recorded, since
a node handed in from outside the base cell must have that wrap added back to
every offset reported for it.

A periodic Delaunay triangulation of $\{\boldsymbol{\xi}_i\}$ on the torus
supplies the connectivity. Its output is expressed as a list of triangles whose
corners are pairs of a node index and a periodic offset,
\begin{equation}
  \text{Corner} = (\,i,\; o_1,\; o_2\,)
  \quad\longmapsto\quad
  \mathbf{X} = \mathbf{x}_i + o_1\mathbf{p}_1 + o_2\mathbf{p}_2 ,
  \label{eq:corner}
\end{equation}
so a triangle may span the period boundary without any node being duplicated in
the cloud. Faces are deduplicated by a translation-invariant key, because the
triangulation is maintained internally on a nine-sheeted covering.

Projection must be injective: two nodes sharing an in-plane position modulo the
periods collapse to one vertex, and the construction rejects that case rather
than silently losing a node. The count is checked against the Euler
characteristic of the torus, $V - E + F = 0$ with $E = \tfrac{3}{2}F$, giving
\begin{equation}
  F = 2V = 2n
  \label{eq:euler}
\end{equation}
faces for $n$ nodes --- one node is already enough to tile the torus, with two
triangles.

\subsection*{Orientation}

The faces must be consistently oriented, and this cannot be decided face by face
from the geometry: a periodic Delaunay triangulation of few points emits
triangles whose three corners are exactly collinear in projection, and a
collapsed triangle has no orientation of its own. Consistent orientation is a
combinatorial property of the surface, so it is propagated across shared edges.
Two faces meeting along a torus edge must traverse it in opposite directions;
where they do not, one is flipped. The edge descriptor is translation
invariant, comprising the two node indices together with the offset difference
between the endpoints,
\begin{equation}
  \text{Edge}(a,b) = \bigl(\, i_a,\, i_b,\; o_1^{b}-o_1^{a},\; o_2^{b}-o_2^{a} \,\bigr).
  \label{eq:edgekey}
\end{equation}
Every edge must be shared by exactly two faces and the propagation must reach
every face, or the connectivity does not describe a surface.

Only once the surface is consistently oriented is the geometry consulted, once,
to fix the global sense: a Delaunay triangle projects to positive signed area
when it faces along $\nhat$, so the largest face --- the best conditioned ---
decides whether the whole surface is reversed.

Finally each face is re-anchored so that its centroid falls in the base period
cell. The deduplication key is translation invariant and so may return any
image as representative; anchoring instead on the lowest node index yields
representatives that overlap in places and leave gaps in others, which corrupts
the ray-crossing count used later for the side test. Anchoring on the centroid
makes the face set a genuine fundamental domain.

\subsection*{Shared connectivity and opposite senses}

The triangulation is evaluated on the \emph{deformed} nodes
$\mathbf{x}_i + \mathbf{u}_i$, and the resulting connectivity is handed to
\emph{both} facets. Since by~\eqref{eq:coincidence} the deformed nodes of the
two grains coincide, sharing the connectivity makes the two deformed surfaces
identical triangle for triangle, that is, one continuous surface. Triangulating
the two reference clouds independently does not achieve this: they differ by the
translations $\mathbf{t}_i$, so independent triangulations produce different
connectivity and the deformed surfaces then agree only at the nodes.

The two facets nevertheless face opposite ways. Each grain must lie on the
positive side of its own facet, otherwise the solid angle --- and with it the
sign of the displacement --- comes out backwards for that grain. This is carried
by a sense $\varsigma \in \{+1,-1\}$ applied to the face normals,
\begin{equation}
  \mathbf{n}_f
  = \varsigma\,
    \frac{(\mathbf{X}_1 - \mathbf{X}_0)\times(\mathbf{X}_2 - \mathbf{X}_0)}
         {\lVert (\mathbf{X}_1 - \mathbf{X}_0)\times(\mathbf{X}_2 - \mathbf{X}_0) \rVert},
  \qquad
  \varsigma_{\A} = -1,
  \quad
  \varsigma_{\B} = +1 .
  \label{eq:sense}
\end{equation}

\subsection*{Realising the mesh}

Every corner~\eqref{eq:corner} referenced by the connectivity becomes a mesh
vertex, at
\begin{equation}
  \mathbf{X}_{(i,o_1,o_2)} = \mathbf{x}_i + o_1\mathbf{p}_1 + o_2\mathbf{p}_2 ,
  \qquad
  \mathbf{U}_{(i,o_1,o_2)} = \mathbf{u}_i ,
  \label{eq:mesh-vertex}
\end{equation}
the displacement of an image being that of its node, since the image
$\mathbf{x}_i + \mathbf{T}$ deforms to
$(\mathbf{x}_i+\mathbf{u}_i) + \mathbf{T}$. The enumeration order depends only
on the connectivity, so two facets sharing it receive identical vertex numbering
and an identical face matrix --- which is what allows them to be compared
directly.

A face may be collapsed in the reference configuration even though it is well
shaped in the deformed one, because the reference facet is stepped: three nodes
spanning a proper triangle after deformation may be exactly collinear before it.
Such a face has no reference area, so its normal is set to zero, which zeroes its
weight in every surface integral, and it is retained in the face matrix so that
both facets keep identical connectivity.

\section*{The plastic displacement field}

\subsection*{The integral}

The displacement at a point $\mathbf{x}$ off the surface is the nodal
displacement weighted by the solid angle the surface subtends there,
\begin{equation}
  \mathbf{u}(\mathbf{x})
  = \frac{1}{2\pi}\int_{\Sinf}
    \mathbf{u}_S(\mathbf{y})\,
    \frac{(\mathbf{x}-\mathbf{y})\cdot\nhat(\mathbf{y})}
         {\lVert \mathbf{x}-\mathbf{y} \rVert^{3}}\,
    \mathrm{d}S(\mathbf{y}) ,
  \label{eq:field}
\end{equation}
where $\Sinf$ is the facet \emph{together with all of its periodic images}. The
images are not a refinement: a single copy of a finite patch subtends a solid
angle that decays with distance instead of tending to the $2\pi$ of a sheet, so
without them the field neither is periodic nor converges to the right limit.

The normalisation $1/2\pi$ is fixed by the case the construction must reproduce.
A doubly periodic sheet subtends exactly $\pm 2\pi$ at any point off it,
whatever its shape, so a uniform nodal displacement $\mathbf{b}$ must give the
rigid translation $\pm\mathbf{b}$ --- which is precisely what a grain translated
by $\pm\tfrac{1}{2}\mathbf{t}$ does.

\subsection*{Discretisation}

The integral is evaluated face by face in closed form. For a triangle with
corners $\mathbf{a},\mathbf{b},\mathbf{c}$ measured from the field point, the
van~Oosterom--Strackee construction gives the signed solid angle
\begin{equation}
  \tan\frac{\Omega}{2}
  = \frac{\mathbf{a}\cdot(\mathbf{b}\times\mathbf{c})}
         {\lVert\mathbf{a}\rVert\lVert\mathbf{b}\rVert\lVert\mathbf{c}\rVert
          + (\mathbf{a}\cdot\mathbf{b})\lVert\mathbf{c}\rVert
          + (\mathbf{a}\cdot\mathbf{c})\lVert\mathbf{b}\rVert
          + (\mathbf{b}\cdot\mathbf{c})\lVert\mathbf{a}\rVert} ,
  \label{eq:vos}
\end{equation}
evaluated with the two-argument arctangent, which is required rather than
convenient: the denominator turns negative once the triangle subtends more than
a hemisphere, and the signed numerator carries the orientation.

A closed form is used because the kernel of~\eqref{eq:field} varies as
$\lVert \mathbf{x}-\mathbf{y}\rVert^{-3}$ and no fixed quadrature rule resolves
it close to the surface; the closed form is exact at every distance, and
cheaper. A collapsed face gives a vanishing numerator and therefore contributes
nothing, as a face with no area should.

\subsection*{Subdivision and the interpolated displacement}

The triangulation joins the nodes, so its triangles are as large as the node
spacing of the boundary. Holding the displacement constant over such a triangle
resolves nothing of how the displacement varies from one node to the next, so
each face is subdivided for the purpose of the quadrature. Cutting every edge
into $k$ parts splits a face into $k^{2}$ similar sub-triangles, whose corners
are the barycentric lattice points
\begin{equation}
  \mathbf{Q}_{ij}
  = \mathbf{X}_0
    + \frac{i}{k}\left( \mathbf{X}_1 - \mathbf{X}_0 \right)
    + \frac{j}{k}\left( \mathbf{X}_2 - \mathbf{X}_0 \right),
  \qquad i,j \ge 0, \quad i+j \le k ,
  \label{eq:subdivision}
\end{equation}
assembled into $\tfrac{1}{2}k(k+1)$ upward triangles
$(\mathbf{Q}_{ij},\mathbf{Q}_{i+1,j},\mathbf{Q}_{i,j+1})$ and
$\tfrac{1}{2}k(k-1)$ downward ones
$(\mathbf{Q}_{i+1,j},\mathbf{Q}_{i+1,j+1},\mathbf{Q}_{i,j+1})$, which is $k^{2}$
in all. The displacement is interpolated linearly across the face by the same
barycentric coordinates, and each sub-triangle carries its value at its own
centroid,
\begin{equation}
  \ubar_{P}
  = \tfrac{1}{3}\left( \mathbf{u}(\mathbf{Q}_{a}) + \mathbf{u}(\mathbf{Q}_{b})
                     + \mathbf{u}(\mathbf{Q}_{c}) \right),
  \label{eq:panel-displacement}
\end{equation}
$\mathbf{Q}_a,\mathbf{Q}_b,\mathbf{Q}_c$ being that sub-triangle's corners.
Because the interpolant is linear, its value at the centroid is also its exact
mean over the sub-triangle: \eqref{eq:panel-displacement} is a midpoint rule,
second order in the panel size, and not merely a point sample. The
displacement remains piecewise constant, but now over panels far smaller than
the node spacing, and it stays bounded whatever the refinement because
$\lvert\Omega\rvert \le 2\pi$.

Subdividing refines only the quadrature. The faces are planar, so
\eqref{eq:subdivision} introduces no new geometry; the side test, the
surface-proximity test and the exported mesh all continue to use the
unsubdivided faces. Two properties of the unrefined scheme survive exactly, and
they are why the refinement is done by subdivision rather than by a
higher-order panel formula:
\begin{itemize}
  \item the sub-triangles tile their parent, so their solid angles sum to the
        parent's and the total solid angle --- and with it the closure
        of~\eqref{eq:closure} --- is unchanged at every refinement;
  \item a uniform nodal displacement interpolates to the same constant on every
        sub-triangle, so the rigid-translation result~\eqref{eq:uniform-check}
        holds for any $k$.
\end{itemize}
Setting $k=1$ recovers the one-value-per-face rule, with
$\ubar_P = \tfrac{1}{3}(\mathbf{u}_{f,0}+\mathbf{u}_{f,1}+\mathbf{u}_{f,2})$.

\subsection*{Periodic images and the far-field closure}

Translating an image of the facet by $\mathbf{T}$ is the same as translating the
field point by $-\mathbf{T}$, so the image sum is evaluated by re-evaluating the
base facet at shifted points. Writing $\Omega_P(\cdot)$ for the solid angle of
panel $P$ and summing over $\lvert m\rvert, \lvert n\rvert \le N$ shells,
\begin{equation}
  \mathbf{W}(\mathbf{x})
  = \sum_{m,n}\sum_P \Omega_P(\mathbf{x}-\mathbf{T}_{mn})\, \ubar_P ,
  \qquad
  \Omega_{\text{expl}}(\mathbf{x})
  = \sum_{m,n}\sum_P \Omega_P(\mathbf{x}-\mathbf{T}_{mn}) .
  \label{eq:shells}
\end{equation}
The subdivision~\eqref{eq:subdivision} is applied to the nearest $3\times3$
block of images only. Beyond that the kernel varies little across a whole face,
so the face's mean displacement is already the correct weight for it and
subdividing would buy nothing; restricting the refinement this way keeps the
cost close to that of the unrefined sum instead of multiplying the entire image
sum by $k^{2}$. Neither the total solid angle nor the rigid-translation case
depends on where that split falls, by the two properties established above.
The explicit shells miss the remainder of the infinite sheet, and that tail
converges only as $1/N$, so truncation alone is not acceptable. Beyond the
explicit shells, however, the kernel varies slowly across a period cell, so
those images see only the area-weighted mean nodal displacement
\begin{equation}
  \ubar = \frac{\sum_f A_f\, \ubar_f}{\sum_f A_f} ,
  \label{eq:mean-displacement}
\end{equation}
while the total solid angle of the whole sheet is known exactly. Replacing the
tail by $\ubar$ times the solid angle still unaccounted for,
\begin{equation}
  \boxed{\;
  \mathbf{u}(\mathbf{x})
  = \frac{1}{2\pi}
    \Bigl[\, \mathbf{W}(\mathbf{x})
      + \ubar\,\bigl( \Omega_{\Sinf}(\mathbf{x})
                    - \Omega_{\text{expl}}(\mathbf{x}) \bigr) \Bigr] ,
  \qquad
  \Omega_{\Sinf}(\mathbf{x}) = 2\pi\,\sigma(\mathbf{x}) ,
  \;}
  \label{eq:closure}
\end{equation}
with $\sigma(\mathbf{x}) = \pm 1$ the side of the facet on which $\mathbf{x}$
lies. For a uniform nodal displacement $\mathbf{u}_i \equiv \mathbf{b}$ we have
$\ubar_f = \ubar = \mathbf{b}$, so~\eqref{eq:closure} collapses to
\begin{equation}
  \mathbf{u}(\mathbf{x})
  = \frac{1}{2\pi}\Bigl[ \mathbf{b}\,\Omega_{\text{expl}}
      + \mathbf{b}\,(2\pi\sigma - \Omega_{\text{expl}}) \Bigr]
  = \sigma\,\mathbf{b}
  \label{eq:uniform-check}
\end{equation}
\emph{exactly}, for any number of shells. The closure therefore absorbs the mean
exactly, and $N$ controls only how well the \emph{variation} about that mean is
resolved far from the facet --- a contribution that converges as $1/N$, since a
distant image contributes only its dipole moment, at a cost growing as
$(2N+1)^2$.

Note that $\Omega_{\Sinf}$ in~\eqref{eq:closure} is obtained from the side test,
not from the truncated series. Summing the series to decide the sign is not
safe: for a point far from the facet, or one lying down inside a corrugation,
the explicit shells can carry the wrong sign entirely.

\section*{Atoms on the boundary, and away from it}

\subsection*{On a node: the field is not evaluated}

If $\mathbf{x}$ coincides with a node of the facet --- allowing for whole
in-plane periods,
\begin{equation}
  \bigl\lVert\, (\mathbf{x} - \mathbf{X}_v)
    - \bigl[\mathbf{p}_1 \mid \mathbf{p}_2\bigr]
      \bigl\lfloor \boldsymbol{\alpha} \bigr\rceil \,\bigr\rVert
  < \varepsilon_{\text{node}},
  \qquad
  \boldsymbol{\alpha}
  = \arg\min_{\boldsymbol{\alpha}}
    \bigl\lVert [\mathbf{p}_1\!\mid\!\mathbf{p}_2]\boldsymbol{\alpha}
      - (\mathbf{x}-\mathbf{X}_v) \bigr\rVert ,
  \label{eq:node-match}
\end{equation}
with $\lfloor\cdot\rceil$ denoting rounding to the nearest integer --- then that
node's own displacement is returned:
\begin{equation}
  \mathbf{u}(\mathbf{x}) = \mathbf{U}_v .
  \label{eq:node-shortcircuit}
\end{equation}
This is not an optimisation. The kernel of~\eqref{eq:field} is singular on the
surface and the solid angle jumps from $+2\pi$ to $-2\pi$ across it, so the
field there is genuinely undefined rather than merely inaccurate. The nodes are
exactly the atoms the boundary brings into coincidence, and their displacements
are prescribed by~\eqref{eq:nodal-displacements}; the quadrature is not asked
for them. Between the nodes the kernel is still singular on the surface and the
discretisation does not resolve it, so accuracy degrades for a point lying on the
facet that is not one of its nodes.

\subsection*{Away from the surface: folding, then the side test}

A point away from the nodes is first folded over the base cell,
\begin{equation}
  \mathbf{x} \leftarrow \mathbf{x}
    - \bigl[\mathbf{p}_1 \mid \mathbf{p}_2\bigr]
      \bigl\lfloor \boldsymbol{\alpha}(\mathbf{x}-\mathbf{x}_c) \bigr\rceil ,
  \label{eq:fold}
\end{equation}
$\mathbf{x}_c$ being the vertex centroid. The field is periodic, so this is
exact; it is also what makes the \emph{truncated} field periodic, since a window
of images fixed in space would otherwise include a different set of images
depending on where the query point sits.

The side $\sigma(\mathbf{x})$ is decided by ray casting: a ray leaving
$\mathbf{x}$ along the facet normal crosses the surface an even number of times
on the positive side and an odd number on the negative side. A nearest-point
test is not sufficient --- for a point down inside a corrugation the nearest
point can lie on the flank of a neighbouring ridge, whose normal faces away, and
the test then reports the wrong side. Three rays tilted slightly apart are cast
and the majority taken, because a ray grazing an edge shared by two faces is
counted twice and yields even parity on either side; with few nodes per cell the
shared edge of two large triangles runs right across it, so this is a systematic
rather than a rare failure, and a degenerate hit cannot catch all three
directions at once.

Both the side test and the surface-proximity test run against a triangle soup
covering a $3\times3$ block of period images, held in a bounding-volume
hierarchy. A soup is used rather than a connected mesh because the faces of a
cut-open periodic surface do not in general glue into a manifold patch: each
face is represented by whichever periodic copy keeps it compact, and
neighbouring faces need not agree on which copy of a shared vertex they use.
Closest-point and ray queries require no connectivity.

\subsection*{Assigning atoms to grains}

Both lattices fill the whole simulation box, so each must be cut back to the
region its grain actually occupies. With $\sigma_{\A}, \sigma_{\B}$ the side
tests of the two facets and $\varsigma_{\A}, \varsigma_{\B}$ the senses
of~\eqref{eq:sense}, the side measured along the outward normal of grain $\A$ is
$\varsigma\,\sigma$, and
\begin{align}
  \mathbf{x} \in \A
  &\iff
  \text{onSurface}_{\A}(\mathbf{x})
  \;\lor\; \varsigma_{\A}\,\sigma_{\A}(\mathbf{x}) < 0 ,
  \label{eq:ingrainA}\\[2pt]
  \mathbf{x} \in \B
  &\iff
  \text{onSurface}_{\B}(\mathbf{x})
  \;\lor\; \varsigma_{\B}\,\sigma_{\B}(\mathbf{x}) > 0 .
  \label{eq:ingrainB}
\end{align}
Grain $\A$ is the material \emph{below} its facet and grain $\B$ the material
\emph{above} its own, which is why the two inequalities run in opposite
directions. Multiplying by $\varsigma$ re-expresses each facet's own answer in
the common frame.

The surface term is not redundant. The side test returns $+1$ or $-1$ and has no
third state, so an atom lying exactly on the surface would be pushed arbitrarily
to one side --- yet those atoms are precisely the nodes, the sites at which the
two grains are brought into coincidence. They belong to the grain rather than to
neither side of it, and~\eqref{eq:node-shortcircuit} then supplies their
displacement directly.

Finally, an atom is displaced by the facet of the grain it belongs to,
\begin{equation}
  \mathbf{u}(\mathbf{x}) =
  \begin{cases}
    \mathbf{u}_{\A}(\mathbf{x}), & \mathbf{x} \in \A ,\\
    \mathbf{u}_{\B}(\mathbf{x}), & \mathbf{x} \in \B ,\\
    \mathbf{0},                  & \text{otherwise},
  \end{cases}
  \label{eq:dispatch}
\end{equation}
each $\mathbf{u}_{(\cdot)}$ being~\eqref{eq:closure} evaluated on that grain's
facet. Atoms of $\A$ above facet $\A$, and atoms of $\B$ below facet $\B$, lie on
the far side of the boundary and are discarded; without that cut the two grains
interpenetrate and the box holds roughly twice the atoms it should.

\section*{Parameters as implemented}

The evaluation carries two discretisation parameters: the subdivision $k$
of~\eqref{eq:subdivision} and the number of explicitly summed image shells $N$
of~\eqref{eq:shells}. Both are set to $4$, so each face is integrated over $16$
panels and $81$ images are summed explicitly, the remainder being closed off
by~\eqref{eq:closure}.

Because the subdivision is confined to the nearest $3\times3$ block, the cost of
one field evaluation is, per face,
\begin{equation}
  9\,k^{2} + \left[ (2N+1)^{2} - 9 \right]
  \;=\; 144 + 72 \;=\; 216
  \label{eq:cost}
\end{equation}
panel evaluations, against $(2N+1)^{2} = 289$ for the unrefined scheme at
$k=1$, $N=8$.

Measured against a converged reference ($k=16$, $N=16$) on a corrugated cell
carrying a spatially varying jump, the discretisation error of the field is
$1.7\times10^{-3}$ at worst and $1.0\times10^{-3}$ on average, in units of the
nodal displacement. The unrefined scheme $k=1$, $N=8$ gives $1.7\times10^{-2}$
and $7.7\times10^{-3}$ on the same test, so the refinement is an order of
magnitude more accurate at rather less cost.

The two parameters are balanced at these values, which is the reason for
choosing them together. At $k=1$ the error is dominated by holding the
displacement constant over node-spacing triangles, and halving $N$ from $8$ to
$4$ changes it by under $4\%$: the far field was never the binding constraint.
Once $k=4$ resolves the near field, raising $N$ from $4$ to $8$ improves the
error by about $1.5$, so neither term is then wasting effort on an error the
other dominates. Greater accuracy is obtained by raising both together; raising
either alone gains little.

The error of the rigid-translation case is unaffected by either parameter, since
\eqref{eq:uniform-check} holds exactly for any $k$ and any $N$.

\section*{Consistency of the two facets}

The construction is only meaningful if the two facets describe one surface.
Three conditions are checked. Corresponding nodes must deform onto the same
point, up to an in-plane period; the two facets must share connectivity; and
their deformed vertices must coincide.

The period tolerance in the first condition is necessary because the two node
sets are wrapped into the simulation box independently, so corresponding nodes
may end up in different periodic images and their deformed positions then agree
only modulo $\mathbf{T}_{mn}$. Before triangulation this is undone: each node of
$\B$ is moved to the image for which
\begin{equation}
  \bigl(\mathbf{x}^{\B}_i + \mathbf{u}^{\B}_i\bigr) - \mathbf{T}
  = \mathbf{x}^{\A}_i + \mathbf{u}^{\A}_i ,
  \qquad
  \mathbf{T} \in \{\mathbf{T}_{mn}\} ,
  \label{eq:glue}
\end{equation}
which costs nothing, since the box vectors are CSL vectors and shifting a node
by one of them lands on another site of the same grain. A residual that is not
an in-plane period is an error, and is reported as such: the two surfaces cannot
then be glued into one.

\end{document}
